\section{Lessons Learned: Famous Algorithms, Tiny Code}
Our main point is that if we take the time to read and refactor
AI code, we find ways to dramatically simplify it. Under the hood,
the \cls{Data}, \cls{Num}, and \cls{Sym} classes turned out to be
so reusable that by the time we had built our active learner, we
already had all the primitives needed for many other inference
tasks---classifiers, clusterers, and more. AI courses typically
devote a lecture each to Na\"ive Bayes, $k$-means, and $k$-means++,
yet once the \cls{Num}/\cls{Sym}/\cls{Data} machinery was in place,
all three collapsed to 30-odd lines between them.

\subsection*{Na\"ive Bayes (in under  30 lines) 
The core of Na\"ive Bayes is one question: how much does a column
{\em like} a value? \textsf{like} answers that. For a \cls{Sym}, it
is a smoothed frequency count from \textsf{has} (the same dictionary
\textsf{add} updates at line~\ref{symaddsub}). For a \cls{Num}, it
is a Gaussian PDF using the running \textsf{mu} and \textsf{sd} that
\textsf{add} maintains via Welford (line~\ref{welford}):
\begin{minted}{python}
the.k=1   # config
def like(col, v, prior):
  """How much a column likes a value."""
  if type(col) == Sym:
    return (col.has.get(v,0) +
      the.bayes.k*prior)/(col.n+the.bayes.k)
  sd=col.sd+1e-32; z=2*sd*sd
  return exp(-(v-col.mu)**2/z)/sqrt(pi*z)
\end{minted}
Note what is {\em not} here: no fitting pass, no separate training
phase. The \cls{Num}/\cls{Sym} columns inside a \cls{Data} are
{\em already} a fitted Bayes model, because \textsf{add} kept them
fitted incrementally as rows arrived. \textsf{like} just reads them.

\textsf{likes} aggregates one row's worth of evidence. It multiplies
per-column likelihoods (in log space, for numerical stability) with
an $m$-estimate prior on class size:
\begin{minted}{python}
the.m = 2 # config

def likes(data, row, n_rows, n_klasses):
  """Log likelihood of row given data."""
  prior = (len(data.rows)+the.bayes.m
           ) / (n_rows+the.bayes.m*n_klasses)
  ls = [like(col,v,prior) for col in data.cols.xs
        if (v:=row[col.at])!="?"]
  return log(prior)+sum(log(v) for v in ls if v>0)
\end{minted}
That is a complete Na\"ive Bayes classifier with Laplace/$m$
smoothing, Gaussian numeric attributes, and missing-value
handling---in two functions. The textbook treatment is a chapter.

To put \textsf{like}/\textsf{likes} to work, \textsf{classify}
implements the standard incremental {\em test-then-train} protocol.
Each row is first classified using the model so far (the row is
assigned to the \cls{Data} with the highest \textsf{likes} score),
then added to its true class's \cls{Data}. Because \textsf{add} is
incremental, this is genuinely streaming---no retraining ever
happens, no batch passes are made:
\begin{minted}{python}
def classify(src, wait=10):
  """Test then train: classify before update."""
  src = iter(src)
  h,cf,all = {}, None, Data([next(src)])
  for n, row in enumerate(src):
    want = row[all.cols.klass.at]
    if n >= wait:
      cf=dinc(want,max(h,key=lambda kl:likes(
        h[kl],row,len(all.rows),len(h))),cf)
    if want not in h: h[want] = clone(all)
    add(all, add(h[want], row))
  return cf
\end{minted}
The \textsf{wait=10} parameter holds off scoring until the model
has seen enough rows to be worth consulting. \textsf{dinc}
accumulates a confusion matrix; \textsf{h} is a dictionary mapping
each class label to its own \cls{Data}; \textsf{all} is the union.
On standard MOOT classification benchmarks (soybean, diabetes,
etc.), \textsf{classify} is competitive with sklearn's batch Na\"ive
Bayes---using two functions of inference logic and one of
bookkeeping.

\subsection*{$k$-means in thirteen lines}

\textsf{kmeans} cycles assignment and update. Assignment uses the
same \textsf{distx} that \textsf{acquireWithCentroid} used. Update
uses the same \textsf{mids} that caches centroids for
\textsf{acquire}. Cluster membership is just a list of \cls{Data}s,
each populated by \textsf{add}:
\begin{minted}{python}
def kmeans(d, rs=None, k=10, n=10, cents=None):
  """Cluster rows into k groups."""
  rs, out = rs or d.rows, []
  cents = cents or choices(rs, k=k)
  for _ in range(n):
    out = [clone(d) for _ in cents]
    for r in rs:
      add(out[min(range(len(cents)),
        key=lambda j:distx(d,cents[j],r))],r)
    cents = [mids(kid) for kid in out if kid.rows]
  return out
\end{minted}
Every primitive in this loop---\textsf{clone}, \textsf{add},
\textsf{distx}, \textsf{mids}---was built for active learning.
$k$-means rides on it, free.

\subsection*{$k$-means++ in eleven lines}

The standard $k$-means++ initialization rule says: pick one centroid
at random, then pick each subsequent centroid with probability
proportional to its squared distance from the nearest existing
centroid. The phrase ``with probability proportional to'' is the
work; everything else is bookkeeping.

EZR already had that work done. \textsf{pick} samples from a
weighted distribution---a \cls{Sym}'s \textsf{has} dictionary, a
\cls{Num}'s Gaussian, or a plain dict of weights---using the same
polymorphism-by-type that \textsf{mid} and \textsf{spread} use:
\begin{minted}{python}
def pick(it, v=None):
  """Sample from distribution."""
  if Sym == type(it): return pick(it.has)
  if Num == type(it):
    tmp=v if v is not None and v!="?" else it.mu
    lo,hi=it.mu-3*it.sd,it.mu+3*it.sd
    new=tmp+it.sd*2*(rand()+rand()+rand()-1.5)
    return lo + (new-lo) % (hi-lo + 1e-32)
  if dict == type(it):
    n = sum(it.values()) * rand()
    for k, v in it.items():
      n -= v
      if n <= 0: return k
\end{minted}
The dict branch is the classic roulette-wheel sampler: accumulate
random mass, return the key where it runs out. \textsf{pick} was
written to support synthetic-data generation in mutation operators,
but the dict branch is exactly what $k$-means++ needs.

So $k$-means++ is just: build the weight dict, sample, repeat:
\begin{minted}{python}
def kpp(d, rs=None, k=10, few=256):
  """k-means++ centroid selection."""
  rs = rs or d.rows
  out = [choice(rs)]
  while len(out) < k:
    t = sample(rs, min(few, len(rs)))
    ws = {i: min(distx(d,t[i],c)**2
                 for c in out)
          for i in range(len(t))}
    out.append(t[pick(ws)])
  return out
\end{minted}
The \textsf{few=256} subsample is a pragmatic shortcut: on large
data sets, scanning all rows for each new centroid is wasteful when
256 random ones give an adequate distribution. The same instinct
that made \textsf{the.few}=128 enough for active learning makes
\textsf{few}=256 enough here. {\em Reading the code} surfaces this
kind of cross-algorithm symmetry; vibe-coding misses it.

\subsection*{The lesson}

A $\sim$30-line section of EZR contains Na\"ive Bayes, $k$-means,
and $k$-means++---three algorithms that fill three textbook chapters
and three weeks of an AI course. They are this short because they
are not really three algorithms. They are three uses of the same
substrate: incrementally maintained \cls{Num}/\cls{Sym}
distributions, summarized via \textsf{mid}/\textsf{spread}, queried
via \textsf{like}/\textsf{distx}, sampled via \textsf{pick}.

This is the kind of structural similarity Dahl, Huang, and Musk
say humans no longer need to see. Without seeing it, you write the
three algorithms three times, in three styles, with three sets of
bugs. With it, you write one substrate and reuse it. The
30-line outcome is not a clever compression. It is what is left
when you read the code carefully enough to notice that the
chapters were the same chapter all along.